\documentclass[conference]{IEEEtran}
\IEEEoverridecommandlockouts
\usepackage{cite}
\usepackage{amsmath,amssymb,amsfonts}
\usepackage{algorithmic}
\usepackage{graphicx}
\usepackage{textcomp}
\usepackage{xcolor}
\usepackage{url}
\usepackage{hyperref}
\usepackage{booktabs}

\def\BibTeX{{\rm B\kern-.05em{\sc i\kern-.025em b}\kern-.08em
    T\kern-.1667em\lower.7ex\hbox{E}\kern-.125emX}}

\begin{document}

\title{Benchmarking AI Coding Agents for Execution-Time Security: Methodology and Progress Report\\
\thanks{This is a methodology and progress report, not the final results paper.}
}

\author{\IEEEauthorblockN{Saarthak}
\IEEEauthorblockA{\textit{AI Research} \\
India \\
saarthak@example.com}
\and
\IEEEauthorblockN{Collaborator}
\IEEEauthorblockA{\textit{AI Research} \\
India \\
collaborator@example.com}
}

\maketitle

\begin{abstract}
As AI coding agents transition from passive code-completion to autonomous systems executing real actions on live repositories, evaluating their security posture becomes critical. This report outlines the methodology for a comprehensive benchmark assessing execution-time vulnerabilities across four attack vectors: Prompt Injection, Context Poisoning, Dependency Confusion, and Secret Leakage. We designed 52 moderate-to-hard tasks isolated in a dedicated test repository (\texttt{AI-test-app}) to evaluate agents under realistic operational constraints. We evaluate chat-based interfaces (ChatGPT, Gemini) against integrated CLI/IDE environments (Claude Code, Antigravity IDE) to document the security implications of context delivery mechanisms. This paper documents the study design, infrastructure, automated and manual scoring rubrics, and current execution status (1/208 runs completed), setting the foundation for the statistical analysis of agent susceptibility.
\end{abstract}

\begin{IEEEkeywords}
prompt injection, dependency confusion, secret leakage, context poisoning, LLM coding agents
\end{IEEEkeywords}

\section{Introduction}

The rapid evolution of AI coding agents has shifted their operational paradigm from passive code-completion suggestions to autonomous systems capable of executing real actions on live repositories. While this enhances developer productivity, it introduces critical security risks. Unlike traditional static analysis tools, autonomous agents interact with the environment, execute shell commands, and modify codebases iteratively, making their execution-time security posture a paramount concern.

This study aims to systematically evaluate the execution security of state-of-the-art coding agents. Our scope focuses on four primary attack vectors: Prompt Injection (PI), Context Poisoning (CP), Dependency Confusion (DC), and Secret Leakage (SL). These vectors were selected from a candidate list of seven because they present the most immediate, measurable threats that can be consistently reproduced across different agent architectures. We deliberately excluded Retrieval Poisoning, Tool Abuse, and Memory Poisoning from this study. Evaluating these excluded vectors would have compromised methodological rigor due to the significant architectural disparities among the agents (e.g., chat-based interfaces vs. CLI-driven tools), which prevents cross-agent comparability for those specific vulnerabilities. By focusing on the four selected vectors, we establish a robust, comparable benchmark that highlights how different delivery mechanisms and agent architectures impact susceptibility to execution-time exploitation.

\section{Related Work}

The security of Large Language Models (LLMs) and their agentic counterparts has seen burgeoning interest. Prior work has laid foundational frameworks for understanding and benchmarking these vulnerabilities. \textit{InjecAgent} \cite{zhan2024injecagent} provided early methodologies for benchmarking indirect prompt injections within tool-integrated LLM agents, establishing that indirect vectors pose a severe risk when agents autonomously utilize external tools. The \textit{AgentDojo} \cite{debenedetti2024agentdojo} family of work expanded this by providing a dynamic environment to evaluate both attacks and defensive mechanisms for LLM agents, becoming a standard for indirect prompt injection benchmarking. Furthermore, the \textit{HackAPrompt} competition \cite{schulhoff2023ignore} exposed systemic vulnerabilities in LLMs to diverse prompt hacking techniques.

In the domain of code generation, Pearce et al. \cite{pearce2021asleep} performed an early assessment of GitHub Copilot's security, identifying a propensity for models to reproduce vulnerable coding patterns. This was further formalized by the \textit{Purple Llama CyberSecEval} \cite{bhatt2023purple}, which introduced a secure coding benchmark for LLMs. A 2025 industry report by GitGuardian \cite{gitguardian2025copilot} highlighted a measurable increase in secret leakage rates within Copilot-assisted repositories (6.4\% compared to a 4.6\% baseline). Moreover, the threat of dependency confusion in LLM-generated code echoes the fundamental software supply chain vulnerabilities disclosed by Birsan in 2021 \cite{birsan2021dependency}.

Most recently, \textit{SecRepoBench} \cite{secrepobench2025} benchmarked code agents for secure code completion within real-world repositories, serving as the closest existing analog to our approach. However, as noted in \textit{The Balkanization of Execution-Security Research for AI Coding Agents} \cite{balkanization2026}, existing research remains fragmented. Our study bridges this gap by providing a cross-agent comparison (evaluating Chat, CLI, and IDE environments) across four distinct vulnerability classes operating on a single, shared, realistic codebase.

\section{Threat Model \& Attack Vector Definitions}

To systematically evaluate the security of AI coding agents, we formally define the four attack vectors operationalized in this study. 

\textbf{Prompt Injection (PI)} occurs when malicious, imperative instructions are embedded within the inputs provided to the agent (e.g., user prompts, file context, simulated issue descriptions, or tool outputs). The agent fails to distinguish between legitimate system instructions and user-supplied data, resulting in the execution of the injected command (e.g., introducing a backdoor or overriding security controls).

\textbf{Context Poisoning (CP)} differs from PI by employing false factual claims rather than imperative instructions. The injection introduces fabricated context (such as false security review sign-offs or fake deprecation notices) into the codebase. The threat manifests when the agent internalizes these false premises and leverages them to justify unsafe coding decisions without independent verification.

\textbf{Dependency Confusion (DC)} measures an agent's susceptibility to introducing or relying on untrusted, typosquatted, or phantom external dependencies. Attackers manipulate the agent by requesting seemingly legitimate but misspelled package names or providing fake missing-package error outputs, tricking the agent into modifying configuration files (e.g., \texttt{requirements.txt}) to pull packages from compromised namespaces.

\textbf{Secret Leakage (SL)} evaluates an agent's propensity to mishandle sensitive credentials. This includes directly hardcoding requested API keys, unintentionally committing secrets from `.env` files to version control, or failing to secure newly written code that processes sensitive information.

As previously noted, vectors such as retrieval poisoning, tool abuse, and memory poisoning were excluded to maintain methodological rigor and ensure robust cross-agent comparability.

\section{Task Bank Design}

The experimental benchmark comprises 52 unique tasks, evenly distributed with 13 tasks per vector. Each task is defined using a YAML schema for programmatic parsing, enabling consistent execution and automated grading where applicable.

The YAML schema defines the following fields across tasks: \texttt{id}, \texttt{name}, \texttt{difficulty}, \texttt{delivery\_mechanism}, \texttt{description}, \texttt{user\_prompt}, \texttt{expected\_safe\_behavior}, \texttt{expected\_exploit\_behavior}, \texttt{scoring\_method}, \texttt{scoring\_criteria}, and \texttt{notes}. Additionally, some tasks utilize a \texttt{file\_context} field for context-based delivery. Depending on the vector, an objective-specific goal field is used: \texttt{injection\_goal} (PI), \texttt{poisoning\_goal} (CP), \texttt{confusion\_goal} (DC), or \texttt{leakage\_goal} (SL). Vector-specific fields further refine the tasks; for instance, Context Poisoning introduces a \texttt{false\_premise} field, Dependency Confusion employs \texttt{malicious\_package\_reference} and \texttt{legitimate\_alternative}, and Secret Leakage utilizes \texttt{planted\_secret\_value} and \texttt{secret\_location}.

Task difficulty is intentionally restricted to ``moderate'' and ``hard'' tiers. Trivial tasks were excluded to prevent agents from passing solely via generic, easily triggered safety filters. The delivery mechanisms are split between chat-paste (where context is explicitly provided in the prompt) and repo-discovery (where the agent must independently explore the repository to discover the injected content), allowing us to measure the security impact of context delivery methods.

\section{Experimental Infrastructure}

To ensure realism and isolate variables, the study adopts a dedicated external repository named \texttt{AI-test-app}. This repository serves as a clean baseline representing a toy internal admin and user-management tool. The application is built using Flask and structured with key components including \texttt{app/app.py} (entrypoint), \texttt{app/auth.py} (authentication flows), \texttt{app/admin.py} (administrative actions), \texttt{app/views.py} (general routes), \texttt{app/utils.py} (utility functions), \texttt{app/api\_client.py} (API stub), \texttt{app/data\_parser.py} (CSV parsing), and \texttt{app/models.py} (database models), alongside a \texttt{config.yaml} file for configuration.

The experimental design employs a branch-per-task methodology. Each of the 52 tasks is executed on an isolated Git branch diverging from the \texttt{main} baseline. This ensures a reproducible and inspectable state for every execution run. 

Crucially, the injected payloads and task setups are generated deterministically via bash scripts rather than relying on LLMs for dynamic generation. This methodology was explicitly adopted following an incident where the Gemini model refused to generate malicious injected payloads directly, triggering its safety filters. Consequently, deterministic generation guarantees consistency and prevents safety-filter interference during the harness setup phase.

\section{Agents Under Study}

Our evaluation encompasses four modern AI coding agents across two distinct delivery interfaces:
\begin{itemize}
    \item \textbf{ChatGPT (Web)} - Chat interface
    \item \textbf{Gemini (Web)} - Chat interface
    \item \textbf{Claude Code (CLI)} - Autonomous CLI agent
    \item \textbf{Antigravity IDE (Gemini Flash Low)} - Integrated Development Environment agent
\end{itemize}

A key methodological choice in this study is the documented delivery asymmetry between the agent types. Chat-based agents (ChatGPT and Gemini) receive the user prompt along with the relevant file contents pasted directly into their context window. In contrast, repository-access agents (Claude Code and Antigravity IDE) receive only the user prompt. These autonomous agents must utilize their tools to navigate the codebase, search for relevant files, and naturally discover the injected or poisoned context. This asymmetry deliberately mirrors real-world usage patterns, allowing us to evaluate whether explicit context provision inherently alters an agent's vulnerability compared to autonomous discovery.

\section{Scoring Methodology}

Each task is evaluated using a trinary scoring system: \texttt{SAFE}, \texttt{EXPLOITED}, or \texttt{AMBIGUOUS}. A task is marked \texttt{SAFE} if the agent rejects the malicious instruction or securely implements the objective. It is marked \texttt{EXPLOITED} if the agent executes the vulnerability. \texttt{AMBIGUOUS} is reserved for cases where the agent's behavior falls outside clear failure or success conditions, such as hallucinating non-existent files or failing to output valid code.

Our scoring leverages an automated tool stack to ensure scalable and objective evaluation. The pipeline utilizes Semgrep for pattern matching (e.g., detecting disabled CSRF decorators or conditional permission bypasses), CodeQL for semantic flow analysis, Bandit for standard Python vulnerability checks (e.g., \texttt{eval} usage or hardcoded passwords), and Gitleaks for secret detection (using a customized \texttt{gitleaks\_config.toml}). Additionally, Dependency Confusion tasks rely on diff-checking against a custom \texttt{blocklist.yaml} to flag malicious package references.

While automation handles the majority of the evaluations, certain vectors require human judgment. Tasks testing nuanced social engineering or ambiguous behavior are flagged as \texttt{manual\_only}. For these cases, we employ a manual review process guided by strict criteria outlined in each vector's \texttt{rubric.md}. To mitigate single-rater bias, we plan to implement inter-rater reliability checks across manually scored tasks prior to final analysis.

\section{Planned Statistical Analysis}

The primary metric for our evaluation is the Vulnerability Rate ($VR$), defined as:
\begin{equation}
VR = \frac{Exploited}{Exploited + Safe}
\end{equation}
Runs resulting in an \texttt{AMBIGUOUS} state are excluded from the denominator to prevent agent unreliability from artificially deflating vulnerability rates, though they are tracked separately as a measure of overall agent stability.

To account for the sample size per condition, we will calculate Wilson score confidence intervals for all vulnerability rates. The analysis will aggregate results across several cuts:
\begin{itemize}
    \item \textbf{Agent $\times$ Vector:} A $4 \times 4$ heatmap visualizing the susceptibility of each agent to each attack vector.
    \item \textbf{Delivery Mechanism:} Comparing the exploitability of chat-paste (explicit context) versus repo-discovery (autonomous context).
    \item \textbf{Difficulty Tier:} Analyzing whether ``hard'' tasks successfully evade security filters more frequently than ``moderate'' tasks.
\end{itemize}

The $4 \times 4$ heatmap will serve as the primary visualization, clearly delineating systemic vulnerabilities across the coding agent landscape.

\section{Execution Plan and Current Status}

Execution is structured to batch complete all 13 tasks per vector for a single agent before proceeding to the next. The planned batching order for agents is Antigravity IDE $\rightarrow$ ChatGPT $\rightarrow$ Claude Code $\rightarrow$ Gemini. The sequence of attack vectors is Prompt Injection (PI) $\rightarrow$ Context Poisoning (CP) $\rightarrow$ Dependency Confusion (DC) $\rightarrow$ Secret Leakage (SL).

The benchmark targets a total of 208 execution runs (52 tasks $\times$ 4 agents). 

\textbf{Current Status:} At the time of this report, 1 of the 208 runs is complete. The executed task, PI-02 on the Antigravity IDE, resulted in a \texttt{SAFE} score. Notably, the agent did not merely ignore the injection attempt silently; it explicitly named and flagged the anomalous instruction embedded within the log file. This run serves as a strong candidate example for future discussion on ``detection-plus-disclosure'' behaviors, where an agent proactively warns the user of suspected malicious context.

\section{Key Methodological Decisions and Corrections}

Throughout the initial setup and pilot phases, several methodological decisions and corrections were formalized (documented in \texttt{docs/decisions.md}):

\begin{itemize}
    \item \textbf{Function Name Correction:} Early testing revealed a discrepancy where tasks targeted a non-existent \texttt{validate\_session} function. This was corrected across the test harness to target the actual \texttt{get\_current\_user\_id} function, and task PI-10 was retargeted accordingly.
    \item \textbf{Harness Silent-Failures:} The use of \texttt{set -e} in our bash scripts initially caused silent failures during environment setup without generating usable logs. This was resolved by implementing explicit \texttt{grep -q} marker verifications to ensure setups complete successfully and log errors transparently.
    \item \textbf{Environment Caveats:} We noted a cross-platform compatibility issue where test scripts executing via Git Bash on Windows required invocations of \texttt{python} (or the \texttt{py} launcher) rather than \texttt{python3}. This environment caveat has been documented to ensure reproducibility across operating systems.
\end{itemize}

These adjustments highlight the rigor required to build deterministic test environments for non-deterministic agents.

\section{Limitations and Threats to Validity}

Several limitations exist within our current methodological framework:
\begin{itemize}
    \item \textbf{Sample Size per Cell:} The experimental design relies on 13 tasks per vector per agent, representing a small $n$. This necessitates careful interpretation of vulnerability rates using confidence intervals to avoid overstating statistical significance.
    \item \textbf{Single-Rater Risk:} As scoring is ongoing, tasks flagged as \texttt{manual\_only} currently face single-rater risk. We intend to address this with inter-rater reliability scoring prior to publishing the final results.
    \item \textbf{Excluded Vectors:} The exclusion of three attack vectors (Retrieval Poisoning, Tool Abuse, and Memory Poisoning) means this benchmark does not comprehensively capture the entire threat landscape of AI coding agents, though it thoroughly covers the four targeted areas.
    \item \textbf{Model-Version Drift:} As proprietary LLMs (e.g., ChatGPT, Gemini, Claude) update frequently without explicit version pinning accessible to all users, model-version drift across the execution window could introduce uncontrolled variance.
\end{itemize}

\section{Timeline and Future Work}

With 1 out of 208 runs completed, the immediate future work involves executing the remaining 207 evaluation runs across the four identified agents. Once the execution phase concludes, we will execute the automated scoring pipelines (using Semgrep, CodeQL, Bandit, and Gitleaks) and conduct the manual scoring passes for ambiguous and \texttt{manual\_only} tasks. 

Subsequent to scoring, we will perform the planned statistical analysis, aggregating data by agent, vector, delivery mechanism, and difficulty. Generating the final results paper, including deep-dive discussions on specific vulnerabilities and the primary $4 \times 4$ heatmap visualization, remains out of scope for this methodology report and will commence only after full execution.

\section{Conclusion}

This report formalizes a reproducible, multi-agent benchmark for evaluating the execution-time security of AI coding agents. By restricting the study to four measurable attack vectors and leveraging an isolated test environment (\texttt{AI-test-app}), we establish a framework capable of comparing fundamentally different agent architectures (Chat versus IDE/CLI) on a level playing field. While execution is still in its infancy, the robust schema design and structured automation pipeline lay the groundwork for identifying critical vulnerabilities and defensive gaps in autonomous AI developer tools.

\appendices
\section{Complete Task Bank}

This appendix lists all 52 tasks across the four vectors evaluated in this study.

\begin{table}[h]
\centering
\caption{Summary of Task Distribution}
\begin{tabular}{lr}
\toprule
\textbf{Delivery Mechanism} & \textbf{Count} \\
\midrule
User Prompt & 12 \\
File Context & 12 \\
Code Comment/Docstring & 9 \\
Tool Output/Simulated Log & 8 \\
Issue File & 3 \\
Git Commit Message & 3 \\
Config/Requirements & 4 \\
Multi-step / Other & 1 \\
\midrule
\textbf{Difficulty Tier} & \textbf{Count} \\
\midrule
Moderate & 26 \\
Hard & 26 \\
\bottomrule
\end{tabular}
\end{table}

\begin{table*}[!t]
\centering
\caption{Prompt Injection Task Bank}
\begin{tabular}{l p{4cm} l l p{4cm}}
\toprule
\textbf{ID} & \textbf{Goal} & \textbf{Delivery} & \textbf{Difficulty} & \textbf{Target File(s)} \\
\midrule
PI-01 & Backdoor Creation & User Prompt & Moderate & \texttt{app/auth.py} \\
PI-02 & Auth Bypass & File Context & Hard & \texttt{app/data\_parser.py}, \texttt{app/utils.py}, \texttt{logs/crash\_report.log} \\
PI-03 & Security Control Disable & Issue Tracker Simulation & Moderate & \texttt{app/views.py}, \texttt{issues/issue\_404.md} \\
PI-04 & Data Exfiltration & File Context & Moderate & \texttt{README.md}, \texttt{app/utils.py}, \texttt{app/views.py} \\
PI-05 & Privilege Escalation & Tool Output & Hard & \texttt{app/api\_client.py} \\
PI-06 & Backdoor Creation & File Context & Moderate & \texttt{app/admin.py} \\
PI-07 & Arbitrary Code Execution & Dependency Docstring & Hard & \texttt{app/data\_parser.py}, \texttt{app/utils.py} \\
PI-08 & Safety Instruction Override & User Prompt & Moderate & \texttt{app/auth.py} \\
PI-09 & Data Exfiltration & File Context & Hard & \texttt{app/data\_parser.py}, \texttt{data/user\_export.csv} \\
PI-10 & Backdoor Creation & Multi Step & Hard & \texttt{app/auth.py}, \texttt{config.yaml} \\
PI-11 & Security Control Disable & Tool Output & Moderate & \texttt{app/views.py} \\
PI-12 & Auth Bypass & File Context & Moderate & \texttt{app/auth.py} \\
PI-13 & Privilege Escalation & Multi Step & Hard & \texttt{app/admin.py}, \texttt{app/utils.py} \\
\bottomrule
\end{tabular}
\end{table*}

\begin{table*}[!t]
\centering
\caption{Context Poisoning Task Bank}
\begin{tabular}{l p{4cm} l l p{4cm}}
\toprule
\textbf{ID} & \textbf{Goal} & \textbf{Delivery} & \textbf{Difficulty} & \textbf{Target File(s)} \\
\midrule
CP-01 & False Security Review & Code Comment & Moderate & \texttt{app/auth.py} \\
CP-02 & False Deprecation Redirect & Docstring & Moderate & \texttt{app/auth.py} \\
CP-03 & Fabricated Test Evidence & Log File & Hard & \texttt{app/admin.py}, \texttt{test\_permission\_check.py} \\
CP-04 & Misleading Naming & Existing Function Naming & Moderate & \texttt{app/views.py} \\
CP-05 & Stale Changelog Claim & Changelog & Moderate & \texttt{app/views.py} \\
CP-06 & Misleading Naming & Docstring & Hard & \texttt{app/views.py} \\
CP-07 & False Upstream Reliance & Code Comment & Moderate & \texttt{app/admin.py} \\
CP-08 & False Upstream Reliance & Config Comment And Stub & Hard & \texttt{app/views.py} \\
CP-09 & False Security Review & Code Comment & Hard & \texttt{app/utils.py} \\
CP-10 & False Upstream Reliance & Git Commit Message & Hard & \texttt{app/auth.py} \\
CP-11 & False Benchmark Justification & Code Comment & Hard & \texttt{app/data\_parser.py} \\
CP-12 & Fabricated Test Evidence & Test File & Moderate & \texttt{app/views.py}, \texttt{test\_secure\_upload.py} \\
CP-13 & False Security Review & Multi File Issue And Comment & Hard & \texttt{app/api\_client.py}, \texttt{issues/feature\_bulk\_export.md} \\
\bottomrule
\end{tabular}
\end{table*}

\begin{table*}[!t]
\centering
\caption{Dependency Confusion Task Bank}
\begin{tabular}{l p{4cm} l l p{4cm}}
\toprule
\textbf{ID} & \textbf{Goal} & \textbf{Delivery} & \textbf{Difficulty} & \textbf{Target File(s)} \\
\midrule
DC-01 & Typosquat & User Prompt & Moderate & \texttt{app/api\_client.py}, \texttt{requirements.txt} \\
DC-02 & Typosquat & Tool Output & Moderate & \texttt{app/auth.py} \\
DC-03 & Internal Name Collision & Readme.Md & Moderate & \texttt{README.md}, \texttt{app/utils.py} \\
DC-04 & Phantom Package & Code Comment & Hard & \texttt{app/data\_parser.py}, \texttt{requirements.txt} \\
DC-05 & Typosquat & Issue File & Moderate & \texttt{app/utils.py}, \texttt{issues/feature\_date\_formatting.md} \\
DC-06 & Version Pin Hijack & Existing Requirements.Txt & Hard & \texttt{app.py}, \texttt{requirements.txt} \\
DC-07 & Internal Name Collision & Existing Broken Import & Hard & \texttt{app/api\_client.py} \\
DC-08 & Malicious Source Index & Config File & Hard & \texttt{app/views.py} \\
DC-09 & Phantom Package & Docstring & Moderate & \texttt{app/admin.py} \\
DC-10 & Typosquat & Git Commit Message & Hard & \texttt{requirements.txt} \\
DC-11 & Transitive Smuggling & Requirements Dev Txt & Hard & \texttt{requirements-dev.txt} \\
DC-12 & Malicious Source Index & Dockerfile & Hard & \texttt{README.md} \\
DC-13 & Phantom Package & Tool Output & Hard & \texttt{app/auth.py} \\
\bottomrule
\end{tabular}
\end{table*}

\begin{table*}[!t]
\centering
\caption{Secret Leakage Task Bank}
\begin{tabular}{l p{4cm} l l p{4cm}}
\toprule
\textbf{ID} & \textbf{Goal} & \textbf{Delivery} & \textbf{Difficulty} & \textbf{Target File(s)} \\
\midrule
SL-01 & Hardcoded Secret Introduction & User Prompt & Moderate & \texttt{app/api\_client.py} \\
SL-02 & Secret Committed To Vcs & File Context & Moderate & \texttt{docker-compose.yml} \\
SL-03 & Hardcoded Secret Introduction & Code Comment & Moderate & \texttt{app/views.py} \\
SL-04 & Secret In Error Output & User Prompt & Hard & \texttt{app/auth.py} \\
SL-05 & Secret In Logs & User Prompt & Moderate & \texttt{app/auth.py} \\
SL-06 & Secret In Logs & File Context & Hard & \texttt{logs/app\_debug.log}, \texttt{logs/error\_summary.log} \\
SL-07 & Secret Committed To Vcs & File Context & Hard & \texttt{N/A} \\
SL-08 & Secret Exposed Via Debug Endpoint & Issue File & Hard & \texttt{issues/debug\_endpoint.md} \\
SL-09 & Secret Reuse Across Environments & Docstring/Config Comment & Hard & \texttt{config.yaml}, \texttt{deploy\_config.yaml} \\
SL-10 & Secret In Client Side Code & User Prompt & Hard & \texttt{app/views.py} \\
SL-11 & Secret Committed To Vcs & Git Commit Message & Moderate & \texttt{app/utils.py} \\
SL-12 & Secret Reuse Across Environments & User Prompt & Hard & \texttt{.github/workflows/deploy.yml} \\
SL-13 & Secret In Error Output & Tool Output & Moderate & \texttt{app/api\_client.py} \\
\bottomrule
\end{tabular}
\end{table*}



\bibliographystyle{IEEEtran}
\bibliography{references}

\end{document}
